\section*{Capítulo \ref{ch:g10:synthesis-yield} --- La síntesis: rendimiento y pureza}

\begin{solution}{exo:g10:synthesis-yield:1}
$\eta = 4.2 / 5.0 = 0.84$, un \omterm{def:g10:synthesis-yield:yield}{rendimiento} del \qty{84}{\percent}.
\end{solution}

\begin{solution}{exo:g10:synthesis-yield:2}
Pesar los \omterm{def:g7:chemical-reactions:reactant}{reactivos}; calentar a reflujo; filtrar a vacío el \omterm{def:g10:synthesis-yield:synthesis}{producto
bruto}; recristalizar; medir el punto de fusión.
\end{solution}

\begin{solution}{exo:g10:synthesis-yield:3}
Si entra por abajo, el agua llena toda la camisa antes de salir por
arriba, así que se refrigera toda la longitud del refrigerante. La parte
superior se deja abierta para que el aparato no esté cerrado: calentar
un aparato cerrado acumularía presión y podría hacerlo estallar.
\end{solution}

\begin{solution}{exo:g10:synthesis-yield:4}
El producto tiene que ser muy soluble en el \omterm{def:g6:solutions-and-solubility:solute}{disolvente} caliente y poco
soluble en el \omterm{def:g6:solutions-and-solubility:solute}{disolvente} frío. Las impurezas, presentes en pequeña
cantidad, siguen disueltas en el \omterm{def:g6:solutions-and-solubility:solute}{disolvente} frío y acaban en el \omterm{def:g3:separating-mixtures:filter}{filtrado}.
\end{solution}

\begin{solution}{exo:g10:synthesis-yield:5}
No: funde por debajo de \qty{135}{\celsius} y a lo largo de un intervalo
de varios grados, señal de un sólido impuro. Habría que recristalizarla,
secarla y volver a medir su punto de fusión.
\end{solution}

\begin{solution}{exo:g10:synthesis-yield:6}
$n = 2.00 / 138.0 = \qty{0.0145}{mol}$ de ácido salicílico, el \omterm{def:g10:reaction-progress-table:limiting}{reactivo
limitante}; $m_{\max} = 0.0145 \times 180.0 = \qty{2.61}{g}$;
$\eta = 2.03 / 2.61 = 0.78$, es decir, el \qty{78}{\percent}.
\end{solution}

\begin{solution}{exo:g10:synthesis-yield:7}
La masa pesada incluye agua, así que es mayor que la masa de aspirina:
el \omterm{def:g10:synthesis-yield:yield}{rendimiento} calculado es demasiado alto, aquí incluso imposible, por
encima del \qty{100}{\percent}. El \omterm{def:g10:synthesis-yield:synthesis}{producto bruto} seco todavía contiene
impurezas (ácido salicílico sin reaccionar, subproductos), que también
añaden masa: su \omterm{def:g10:synthesis-yield:yield}{rendimiento} también es demasiado alto con respecto al
\omterm{def:g10:synthesis-yield:yield}{rendimiento} real en aspirina pura.
\end{solution}

\begin{solution}{exo:g10:synthesis-yield:8}
Es mucho más rápida, y deja los cristales mucho más secos, ya que la
succión extrae la mayor parte del líquido; además, los cristales son
fáciles de lavar sobre el embudo y de rascar del papel de filtro plano.
\end{solution}

\begin{solution}{exo:g10:synthesis-yield:9}
Producto: $3.1/5.0 = 0.62$; aspirina: $3.1/5.0 = 0.62$; ácido
salicílico: $2.2/5.0 = 0.44$. El producto recorre lo mismo que la
aspirina y no presenta ninguna mancha de ácido salicílico: es aspirina,
sin ácido salicílico detectable.
\end{solution}

\begin{solution}{exo:g10:synthesis-yield:10}
\omterm{def:g6:solutions-and-solubility:solute}{Disolvente} A: muy soluble en caliente y poco soluble en frío; el
\omterm{def:g6:solutions-and-solubility:solute}{disolvente} B mantiene disuelta la mayor parte del producto incluso en
frío. Con A, \qty{3.0}{g} necesitan unos
$3.0 / 150 = \qty{0.020}{L}$, es decir, \qty{20}{mL} de \omterm{def:g6:solutions-and-solubility:solute}{disolvente}
caliente; una vez frío, como mucho siguen disueltos
$2 \times 0.020 = \qty{0.04}{g}$.
\end{solution}

\begin{solution}{exo:g10:synthesis-yield:11}
Un líquido que hierve necesita espacio por encima: la espuma y las
salpicaduras no deben llegar al refrigerante. Los reguladores de
ebullición dan al vapor pequeños puntos donde formar burbujas, para que
el líquido hierva de forma regular y no a borbotones repentinos.
\end{solution}

\begin{solution}{exo:g10:synthesis-yield:12}
$0.80 \times 0.70 = 0.56$: \qty{56}{\percent}. Tres etapas del
\qty{90}{\percent}: $0.90^3 = 0.73$, \qty{73}{\percent}. Cada etapa
pierde producto (y tiempo y dinero), y las pérdidas se multiplican:
cuantas menos etapas, más producto queda al final.
\end{solution}

\begin{solution}{exo:g10:synthesis-yield:13}
Como mucho $3.3 \times 0.050 = \qty{0.17}{g}$ de aspirina. El agua
helada \omterm{def:g2:mixing-and-dissolving:dissolve}{disuelve} todavía menos aspirina, ya que su \omterm{def:g6:solutions-and-solubility:solubility}{solubilidad} baja con
la temperatura.
\end{solution}

\begin{solution}{exo:g10:synthesis-yield:14}
\begin{enumerate}
  \item $n_0(\text{ácido salicílico}) = 2.76 / 138.0 = \qty{0.0200}{mol}$,
        anhídrido, \qty{0.0500}{mol}: el ácido salicílico es el
        limitante, $x_{\max} = \qty{0.0200}{mol}$, y quedan
        $0.0500 - 0.0200 =
        \qty{0.0300}{mol}$ de anhídrido.
  \item La \omterm{def:g10:synthesis-yield:synthesis}{síntesis} da \qty{0.0200}{mol} de ácido etanoico, y la
        destrucción del exceso, $2 \times 0.0300 = \qty{0.0600}{mol}$:
        \qty{0.0800}{mol} en total.
\end{enumerate}
\end{solution}

\begin{solution}{exo:g10:synthesis-yield:15}
$n(\text{aspirina}) = \num{1.00e6} / 180.0 = \qty{5.56e3}{mol}$. Con un
\omterm{def:g10:synthesis-yield:yield}{rendimiento} del \qty{85}{\percent}, el ácido salicílico necesario es
$\num{5.56e3} / 0.85 = \qty{6.54e3}{mol}$, es decir,
$\num{6.54e3} \times 138.0 = \qty{9.0e5}{g}$, unas \qty{0.90}{t}.
\end{solution}

\begin{solution}{pb:g10:synthesis-yield:1}
\textbf{1.} C: $7 + 4 = 11$ y $9 + 2 = 11$; H: $6 + 6 = 12$ y
$8 + 4 = 12$; O: $3 + 3 = 6$ y $4 + 2 = 6$. Ajustada.

\textbf{2.} Ácido salicílico: $7 \times 12.0 + 6 \times 1.0 + 3 \times 16.0
= \qty{138.0}{g/mol}$; anhídrido: $4 \times 12.0 + 6 \times 1.0 +
3 \times 16.0 = \qty{102.0}{g/mol}$; aspirina: $9 \times 12.0 + 8 \times
1.0 + 4 \times 16.0 = \qty{180.0}{g/mol}$; ácido etanoico: $2 \times 12.0
+ 4 \times 1.0 + 2 \times 16.0 = \qty{60.0}{g/mol}$.

\textbf{3.} \omterm{def:g7:chemical-reactions:reactant}{Reactivos}: el ácido salicílico y el anhídrido etanoico;
producto buscado: la aspirina; subproducto: el ácido etanoico.

\textbf{4.} Calentar acelera la reacción; a reflujo, los vapores se
condensan y vuelven, así que no se pierde ningún \omterm{def:g7:chemical-reactions:reactant}{reactivo} ni producto.

\textbf{5.} $n = 3.0 / 138.0 = \qty{0.0217}{mol}$.

\textbf{6.} $m = 6.0 \times 1.08 = \qty{6.48}{g}$;
$n = 6.48 / 102.0 = \qty{0.0635}{mol}$.

\textbf{7.} Ácido salicílico: $0.0217 - x$; anhídrido: $0.0635 - x$;
aspirina y ácido etanoico: $x$. El ácido salicílico se agota primero:
$x_{\max} = \qty{0.0217}{mol}$, el ácido salicílico es el limitante.

\textbf{8.} $m_{\max} = 0.0217 \times 180.0 = \qty{3.91}{g}$.

\textbf{9.} Para que todo el ácido salicílico se convierta en aspirina:
el que sobrara quedaría \omterm{def:g2:mixing-and-dissolving:mix}{mezclado} con la aspirina y sería difícil de
eliminar, mientras que el anhídrido en exceso se destruye con el agua
que se añade al final, y el ácido etanoico que da se elimina en los
lavados.

\textbf{10.} Incluye agua: \qty{3.6}{g} darían un falso
$3.6 / 3.91 = \qty{92}{\percent}$.

\textbf{11.} $3.3 / 3.91 = 0.84$: un \qty{84}{\percent} de \omterm{def:g10:synthesis-yield:synthesis}{producto
bruto} seco, que no es puro.

\textbf{12.} Se eliminan las impurezas, y parte de la aspirina se queda
disuelta en el \omterm{def:g6:solutions-and-solubility:solute}{disolvente} frío y en las aguas de lavado.

\textbf{13.} $3.3 \times 0.030 = \qty{0.10}{g}$ como mucho.

\textbf{14.} $\eta = 2.7 / 3.91 = 0.69$: \qty{69}{\percent}.

\textbf{15.} Un punto de fusión neto en el valor de la aspirina pura:
los cristales son aspirina, y pura.

\textbf{16.} Funde por debajo de \qty{135}{\celsius} y en un intervalo
amplio: el \omterm{def:g10:synthesis-yield:synthesis}{producto bruto} era impuro.

\textbf{17.} El \omterm{def:g10:synthesis-yield:synthesis}{producto bruto} contenía aspirina y ácido salicílico sin
reaccionar; la \omterm{def:g10:synthesis-yield:recrystallisation}{recristalización} eliminó el ácido salicílico.

\textbf{18.} $R_f = 3.1 / 5.0 = 0.62$.

\textbf{19.} Más o menos lo mismo en cada una: de los \qty{3.91}{g}
posibles a \qty{3.3}{g} de \omterm{def:g10:synthesis-yield:synthesis}{producto bruto}, y de \qty{3.3}{g} a
\qty{2.7}{g} en la \omterm{def:g10:synthesis-yield:recrystallisation}{recristalización}, unos \qty{0.6}{g} cada vez. Como el
\omterm{def:g10:synthesis-yield:synthesis}{producto bruto} no era aspirina pura, la primera parte del trabajo perdió
en realidad algo más de \qty{0.6}{g} de aspirina.

\textbf{20.} El \omterm{def:g10:synthesis-yield:yield}{rendimiento} de la \omterm{def:g10:synthesis-yield:synthesis}{síntesis} es de alrededor del
\qty{69}{\percent}.
\end{solution}
